\begin{abstract}
\noindent\textbf{Аннотация.} В музыкальных стриминговых сервисах до 35\% пользовательских действий представляют собой быстрые скипы ($<30$ секунд), отражающие неудовлетворённость рекомендацией. Стандартная оптимизация правдоподобия (Maximum Likelihood Estimation, MLE) штрафует авторегрессионные трансформеры за неспособность предсказать такие негативные события (проблема \textit{False Negatives}), тиражируя нежелательный контент. В данной работе исследуется выравнивание последовательного рекомендателя на базе причинной архитектуры Causal GPT-2 (43.2M параметров) с целевой удовлетворённостью пользователя посредством прямого согласования предпочтений (Direct Preference Optimization, DPO) с якорной SFT-регуляризацией ($\gamma=0.10$). Эксперименты на масштабном корпусе Яндекс Музыки (\textbf{Yambda-50M}: 98\,505 токенов каталога, 33.2M обучающих взаимодействий, 9\,048 отложенных тестовых пользователей) показали, что DPO не только превосходит 2-стадийные методы реранжирования с оффлайн-критиком и PPO, но и полностью устраняет «налог на выравнивание» (\textit{Zero Alignment Tax}). Политика DPO обеспечивает статистически достоверный прирост точности ранжирования (\textbf{+0.67\% HR@10}, $p < 0.0001$; тест МакНемара $\chi^2 = 18.32, p = 1.9 \times 10^{-5}$), расширяет покрытие каталога на \textbf{+1\,027 уникальных треков} из длинного хвоста (+1.04\% каталога), снижает зависимость от поп-хитов на \textbf{$-12.1\%$} и увеличивает кумулятивное время непрерывного прослушивания в 5-шаговых роллаутах на \textbf{+1.5 секунды} без деградации из-за \textit{Exposure Bias}.

\vspace{0.5em}
\noindent\textbf{Ключевые слова:} последовательные рекомендации, обучение с подкреплением, Direct Preference Optimization, Causal GPT-2, стриминг музыки, проблема скипов, Zero Alignment Tax, Yambda.
\end{abstract}

\vspace{0.5em}
\begin{otherlanguage}{english}
\begin{abstract}
\noindent\textbf{Abstract.} In music streaming platforms, up to 35\% of user interactions correspond to immediate skips ($<30$s), reflecting user dissatisfaction. Standard Maximum Likelihood Estimation (MLE) penalizes causal autoregressive transformers for failing to predict such negative interactions (the \textit{False Negative} dilemma), inadvertently amplifying disliked content. In this work, we investigate the alignment of a Causal GPT-2 recommender (43.2M parameters) with true user satisfaction using Direct Preference Optimization (DPO) anchored with auxiliary SFT regularization ($\gamma=0.10$). Extensive evaluations on the industrial-scale Yandex Music dataset (\textbf{Yambda-50M}: 98,505 catalog tokens, 33.2M training interactions, 9,048 held-out test users) demonstrate that DPO not only outperforms 2-stage Critic Re-Ranking and Actor-Critic PPO, but completely circumvents the alignment tax (\textit{Zero Alignment Tax}). The aligned DPO policy achieves statistically significant next-item ranking gains (\textbf{+0.67\% HR@10}, $p < 0.0001$; McNemar's paired test $\chi^2 = 18.32, p = 1.9 \times 10^{-5}$), expands catalog coverage by \textbf{+1,027 unique tracks} from the long tail, reduces popularity Spearman correlation by \textbf{$-12.1\%$}, and yields a \textbf{+1.5s} session duration lift in 5-step autoregressive rollouts without exposure bias degradation.

\vspace{0.5em}
\noindent\textbf{Keywords:} Sequential Recommendation, Reinforcement Learning Alignment, Direct Preference Optimization, Causal GPT-2, Music Streaming, Skip Dilemma, Zero Alignment Tax, Yambda-50M.
\end{abstract}
\end{otherlanguage}
